\documentclass[11pt]{article}

\usepackage[margin=1in]{geometry}
\usepackage[T1]{fontenc}
\usepackage{lmodern}
\usepackage{microtype}
\usepackage{hyperref}
\usepackage{enumitem}
\usepackage{array}
\usepackage{longtable}
\usepackage{xcolor}

\hypersetup{
  colorlinks=true,
  linkcolor=blue!50!black,
  urlcolor=blue!50!black,
  citecolor=blue!50!black
}

\setlist[itemize]{leftmargin=1.5em}
\setlist[enumerate]{leftmargin=1.5em}

\newcommand{\tool}{\textsc{mprewriter}}
\newcommand{\hittrace}{\texttt{HITTRC01}}
\newcommand{\code}[1]{\texttt{#1}}

\title{\tool: Low-Overhead Probe-Based Observability and Coverage for Distributed Systems}
\author{Jayanta Majumdar \and Sambuddha Majumder}
\date{\today}

\begin{document}

\maketitle

\begin{abstract}
Distributed systems are difficult to test and diagnose because the behavior of
interest is often spread across languages, processes, hosts, and runtime
contexts. Conventional code coverage tools usually report what happened inside
one process, one test run, or one language runtime. Logging and tracing systems
can describe request paths, but they are often too coarse or too expensive to
use as a dense coverage signal. This paper describes \tool, a small
language-neutral observability substrate for distributed systems. Applications
are instrumented with lightweight probes. A probe records a compact hit message
containing an application id, instance id, thread id, stack depth, and location
id. Probe runtimes batch these messages and send them over UDP or TCP to
collector servers. Collectors persist the raw event history in a common binary
trace format, maintain context-aware hit counts, and answer interactive
coverage queries. A prototype implementation includes Java, C++, and Rust
probe runtimes; Java bytecode instrumentation using ASM; Java, C++, and Rust
collectors; trace dumpers in multiple languages; context management; and
regression tests that exercise cross-language producer and server
combinations. The central design claim is that dense observability becomes
practical when the hot path is reduced to a minimal enqueue, while parsing,
aggregation, durability, and query handling are moved to asynchronous
collectors and offline tools.
\end{abstract}

\section{Introduction}

Coverage is normally understood as a property of a test run: a program was
executed, a coverage tool observed which lines or branches ran, and a report was
produced afterward. That model is useful, but it does not fit many distributed
systems. A modern service boundary may be crossed by Java, C++, and Rust
components. A request may trigger asynchronous tasks and callbacks. A feature
flag, tenant, experiment, market data condition, or failure mode may be the
reason a particular path executes. In that setting, the more useful question is
not only ``was this line covered?'' but ``which runtime contexts exercised this
location, in which application instance, and what history led to it?''

\tool\ is a pragmatic answer to that question. Its design is deliberately
small. Instrumented code calls a probe function with a location id. The probe
runtime records a hit in a process-local buffer and a background sender batches
hits into datagrams. A server receives hit, log, and context messages, persists
the raw arrival history, and maintains counters keyed by application,
instance, context set, and location. The same wire format is shared by Java,
C++, and Rust. The same administrative command vocabulary is exposed by the
Java, C++, and Rust collectors.

The name \tool\ comes from ``multipurpose rewriter''. The original emphasis was
on source or bytecode rewriting, especially inserting probes into application
code. The current implementation still contains that instrumentation-first
heritage, but the larger idea is now a cross-language probe protocol and a
family of collectors and analysis tools.

\section{Goals and Non-Goals}

The project is guided by five goals.

\begin{enumerate}
  \item \textbf{Low hot-path overhead.} A probe should be cheap enough to place
  at method entries, branch edges, exception handlers, and manually selected
  application events.

  \item \textbf{Language neutrality.} A Java service, a C++ service, and a Rust
  service should be able to send equivalent hit messages to the same collector.

  \item \textbf{Distributed coverage.} Coverage should be computable across
  process boundaries and runtime instances, not only within a single test JVM
  or single binary.

  \item \textbf{Runtime context slicing.} Coverage should be attributable to
  dynamic context sets such as feature flags, scenarios, experiments, tenants,
  or operational phases.

  \item \textbf{Post-hoc analysis.} Collectors should persist enough raw
  history that later tools can dump, filter, replay, or summarize a run without
  requiring the original application processes.
\end{enumerate}

The project is not intended to replace distributed tracing systems such as
OpenTelemetry, nor does it attempt to provide strong transactional delivery
semantics by default. The common path uses UDP because losing a small number of
coverage hits may be acceptable for high-volume observation. TCP support and
trace persistence are available when stronger collection is required. The
design favors useful, cheap, and portable observability over a heavyweight
exactness guarantee.

\section{Prototype Implementation}

The prototype implementation is composed of several separable units. The
following table describes them by responsibility rather than assuming any
particular source tree layout. Concrete module names are included only as
identifiers for the implementation artifacts.

\begingroup
\small
\begin{longtable}{>{\raggedright\arraybackslash}p{0.36\linewidth}%
                  >{\raggedright\arraybackslash}p{0.57\linewidth}}
\textbf{Implementation artifact} & \textbf{Role} \\
\hline
\code{code-analytics} & Java collector module. It listens on UDP port 8083 and
TCP port 8084, writes \hittrace\ files, maintains hit counters, and exposes a
Clojure-backed administrative shell. \\
\code{code-analytics-cpp} & C++17 collector module. It provides UDP/TCP
listeners, a binary trace writer, context management, and an embedded Tcl
shell. \\
\code{code-analytics-rust/cov\_server} & Rust collector crate. It uses a Molt
shell, shared protocol parsing, context-aware hit counts, and trace hooks. \\
\code{JavaAppWithUdpProbes} & Java demonstration application and early Java
probe runtime. It exercises manual probes, logs, stack depth, and tests around
a small trading-domain application. \\
\code{CppAppWithUdpProbes} & C++ demonstration application and optimized probe
runtime using a batched UDP sender and a lock-free MPSC ring for hit records. \\
\code{RustAppWithUdpProbes} & Rust demonstration applications and probe
runtimes, including a macro-based interface and latency benchmarks. \\
\code{branch-probe-instrumenter} & Standalone ASM-based Java JAR
instrumenter. It injects method, branch, switch, catch, and finally probes and
emits a probe mapping CSV. \\
\code{branch-probe-suite} & Combined Java instrumentation suite containing the
runtime and instrumenter modules intended to be built and used together. \\
\code{tools/trace-dump-*} & Standalone Java, C++, and Rust trace dumpers for
the common \hittrace\ binary trace format. \\
\code{testsuite} & Scripted cross-language regression harness with golden
outputs. \\
\code{benchmark.md} and \code{run\_benchmarks.ps1} & Engineering notes and a
helper script used to record latency measurements for Java, C++, and Rust probe
paths. \\
\end{longtable}
\endgroup

The implementation history shows a steady evolution from a Java UDP hit server and manual
probes to a broader system: context management, binary trace writing, Rust and
C++ producers and servers, stack-depth propagation, a cross-language test
suite, Java JAR instrumentation, inclusion and exclusion filtering, sidecar
probe maps, trace dumpers, and latency-focused runtime optimizations.

\section{System Model}

\subsection{Actors}

\tool\ separates the system into four actors.

\begin{itemize}
  \item \textbf{Instrumenters} insert calls to a probe API. In Java this is
  done by an ASM JAR rewriter. In C++ and Rust, probes can be inserted by
  source-level tooling, macros, or manual calls.

  \item \textbf{Probe runtimes} live inside the application process. They pack
  location hits into compact records, maintain cheap stack-depth information
  where possible, and send records asynchronously.

  \item \textbf{Collectors} receive UDP and TCP messages, write the raw
  history to a trace file, maintain counters, and expose query commands.

  \item \textbf{Offline tools} dump traces, compute statistics, filter by time
  ranges, compare output to golden files, and correlate location ids with
  source maps.
\end{itemize}

\subsection{Data Flow}

The normal data flow is:

\begin{enumerate}
  \item A probe site executes in application code.
  \item The probe runtime records a hit in a local queue or ring buffer.
  \item A background sender batches one or more records into a UDP datagram.
  \item A collector receives the datagram and writes the exact payload into a
  binary trace frame.
  \item The collector parses the payload and increments coverage counters.
  \item An operator or test script queries the collector using the shared admin
  command vocabulary.
\end{enumerate}

This design keeps application threads away from file I/O, network parsing,
aggregation, and interactive query logic. The application pays only for
location-id capture and enqueue work on the hot path.

\section{Wire Protocol}

All multi-byte fields are big-endian. The most important message is the
20-byte hit record:

\begin{center}
\begin{tabular}{lll}
\textbf{Field} & \textbf{Type} & \textbf{Meaning} \\
\hline
\code{messageType} & \code{u16} & \code{1} for hit \\
\code{appId} & \code{u16} & Application family or logical service id \\
\code{instanceId} & \code{u32} & Runtime instance id \\
\code{threadId} & \code{u32} & Producer thread id or stable thread hash \\
\code{stackDepth} & \code{u32} & Probe-observed call depth \\
\code{locationId} & \code{u32} & Instrumented source or bytecode location \\
\end{tabular}
\end{center}

The log record extends the same envelope:

\begin{center}
\begin{tabular}{lll}
\textbf{Field} & \textbf{Type} & \textbf{Meaning} \\
\hline
\code{messageType} & \code{u16} & \code{2} for log \\
\code{appId} & \code{u16} & Application id \\
\code{instanceId} & \code{u32} & Runtime instance id \\
\code{threadId} & \code{u32} & Producer thread id \\
\code{stackDepth} & \code{u32} & Current stack depth \\
\code{msgLen} & \code{u16} & UTF-8 byte length \\
\code{msg} & bytes & UTF-8 payload \\
\end{tabular}
\end{center}

Context updates are smaller:

\begin{center}
\begin{tabular}{lll}
\textbf{Message type} & \textbf{Payload} & \textbf{Meaning} \\
\hline
\code{3} & UTF-8 label & Attach context label to the active context set \\
\code{4} & UTF-8 label & Withdraw context label from the active context set \\
\end{tabular}
\end{center}

The protocol intentionally does not encode language-specific type systems,
stack frame objects, or source names. Those are reconstructed using the probe
map emitted by the instrumenter or by application-specific metadata. The
on-the-wire contract remains small enough for simple implementations in
multiple languages.

\section{Probe Runtime Design}

\subsection{The Hot Path}

The core optimization is to make a hit record cheap to produce. In the fastest
paths, a probe call:

\begin{enumerate}
  \item obtains a thread id and current depth,
  \item packs fields into a primitive value or a preallocated slot,
  \item claims a queue or ring-buffer position,
  \item marks the slot ready,
  \item returns to application code.
\end{enumerate}

Network I/O is handled by a sender thread. Batching uses 72 hit records per
datagram, because \(72 \times 20 = 1440\) bytes, which fits comfortably under a
typical Ethernet MTU after IP and UDP headers. This avoids fragmentation while
amortizing UDP send cost over many hits.

\subsection{Java Runtime}

The Java runtime appears in several forms. One demonstration runtime, packaged
with a small trading-domain application, uses \code{ArrayBlockingQueue} and
\code{StackWalker} to compute application-frame depth. An optimized benchmark
runtime uses a bounded MPSC ring implemented with arrays, \code{AtomicLong}
head and tail counters, and \code{VarHandle} acquire/release operations for
slot readiness. It can compute depth using \code{StackWalker}, or it can use a
constant-depth mode for extremely hot locations.

The Java branch-probe runtime is a small companion library for bytecode-
instrumented applications. It uses a thread-local depth counter maintained by
instrumented method entry and exit calls. This is important for Java because
full stack walking is expensive. By making the instrumenter insert
\code{scope\_ENTER} and \code{scope\_EXIT}, the runtime can report meaningful
depth with a cheap thread-local integer.

\subsection{C++ Runtime}

The C++ probe runtime uses a lock-free MPSC ring with prepacked 20-byte
network-order frames. Each producer claims a slot by advancing an atomic tail,
writes the frame into a cache-aligned slot, and sets a ready flag. The sender
drains the ring into a reusable send buffer and emits UDP datagrams. Logs are
rarer and use a simple mutex-protected queue.

C++ also has a natural way to maintain stack depth. A probe macro creates an
RAII guard. The guard increments a thread-local depth counter on construction,
records the hit, and decrements the counter on destruction. This makes depth
robust to ordinary control flow and exceptions.

\subsection{Rust Runtime}

The Rust runtimes use the same basic design: an application-facing probe API,
thread-local depth, and an asynchronous UDP worker. The Rust macro support
includes a procedural macro crate, \code{mprewriter\_macros}, that expands a
probe site into creation of a depth guard followed by a call to
\code{scope\_START}. As in C++, Rust's ownership and drop semantics provide a
good fit for depth guards.

The Rust server and Rust probe demos demonstrate the portability of the
protocol. A Rust producer can send hits to a Java server, and C++ producers can
be consumed by Rust or Java servers in the regression suite.

\section{Instrumentation}

\subsection{Java Bytecode Instrumentation}

The Java instrumenter uses ASM to rewrite classes inside a JAR. It injects
probes for:

\begin{itemize}
  \item method entry,
  \item conditional branch true and false edges,
  \item switch targets,
  \item catch handler entries,
  \item finally handler entries.
\end{itemize}

For normal methods, method-entry probes are delayed until the first source line
is available. For constructors, entry probes are delayed until after the
superclass constructor has run. This avoids invalid bytecode around constructor
prologues and gives the probe a useful source line. If a method has no line
information, the instrumenter falls back to emitting a method-entry probe near
exit with an unknown line.

Each injected probe receives a monotonically assigned location id. The
instrumenter writes a CSV index with the columns:

\begin{center}
\code{id,class,method,where,source,line}
\end{center}

The CSV is embedded as \code{META-INF/branch-probes.csv}; optionally, a sidecar
CSV is written next to the output JAR. This map is the bridge between compact
runtime location ids and human-readable source locations.

\subsection{Filtering}

Probe density matters. Instrumenting every branch in a large application can
create too much traffic or noise. The instrumenter therefore supports:

\begin{itemize}
  \item \code{bp.inject}, a comma-separated set of probe kinds to inject,
  \item \code{bp.excludefile}, a class and line-range exclusion file,
  \item \code{bp.includefile}, a class and line-range inclusion file.
\end{itemize}

The inclusion and exclusion mechanism lets a user focus on a package, a
suspect class, or a few line ranges while keeping the same downstream protocol.

\section{Collectors}

\subsection{Common Responsibilities}

All collectors implement the same conceptual loop:

\begin{enumerate}
  \item receive UDP datagrams and TCP streams,
  \item identify the message type,
  \item write the raw payload to the trace file with a source tag,
  \item parse complete hit, log, and context messages,
  \item update in-memory coverage state,
  \item answer administrative queries.
\end{enumerate}

The Java server uses a Clojure-backed shell. The C++ server embeds Tcl. The
Rust server uses Molt. Despite the different shell implementations, the
colon-prefixed command set is intentionally aligned:

\begin{center}
\small
\begin{tabular}{>{\raggedright\arraybackslash}p{0.38\linewidth}%
                >{\raggedright\arraybackslash}p{0.48\linewidth}}
\textbf{Command} & \textbf{Purpose} \\
\hline
\code{:help} & List commands \\
\code{:hits} & Print received or aggregated hits \\
\code{:apply-context <label>} & Add a label to the active context set \\
\code{:withdraw-context <label>} & Remove a label from the active context set \\
\code{:coverage-report <appId> <instanceId> <file>} & Write a coverage report \\
\code{:flush-trace} & Flush buffered trace data \\
\code{:trace-persist} & Force durable persistence \\
\code{:exit} & Stop the shell \\
\end{tabular}
\end{center}

\subsection{Context Management}

The collector maintains a global active context set. The empty context set is
assigned id 1 and is reported as \code{default}. Each distinct non-empty set of
labels is canonicalized and assigned a stable integer id. When a hit arrives,
it is counted against the current context-set id.

This is a simple model, but it is powerful enough for many testing and
observability scenarios. A test driver can apply context \code{feature-X},
exercise a workflow, withdraw it, apply \code{feature-Y}, and later ask which
locations were covered under each context. A production diagnostic session can
mark a phase, incident, synthetic transaction, or customer segment without
changing the application probe sites.

\subsection{Coverage Reports}

The common coverage report format is plain text:

\begin{verbatim}
CONTEXTS N
<ctxId> <label>
...
HITS M
<ctxId> <locId> <count>
...
\end{verbatim}

The report is filtered by application id and instance id. The CSV probe map
then resolves \code{locId} into class, method, branch kind, source file, and
line. Separating the report from the source map keeps the runtime protocol
small and lets instrumentation metadata be archived alongside the build
artifact.

\section{Trace Persistence}

Collectors write a binary trace format with magic \hittrace. The header is:

\begin{center}
\begin{tabular}{lll}
\textbf{Field} & \textbf{Type} & \textbf{Meaning} \\
\hline
\code{magic} & 8 bytes & ASCII \code{HITTRC01} \\
\code{endian} & \code{u8} & \code{0} for big-endian \\
\code{startMillis} & \code{u64} & File creation time in epoch milliseconds \\
\end{tabular}
\end{center}

Each frame is:

\begin{center}
\begin{tabular}{lll}
\textbf{Field} & \textbf{Type} & \textbf{Meaning} \\
\hline
\code{flag} & \code{u16} & Hit, log, context attach, context withdraw, or timestamp \\
\code{source} & \code{u8} & UDP, TCP, or internal \\
\code{nanos} & \code{u64} & Collector-side monotonic timestamp \\
\code{len} & \code{u32} & Payload length \\
\code{payload} & bytes & Raw message bytes \\
\end{tabular}
\end{center}

The frame payload is deliberately raw. If a UDP datagram contains multiple
20-byte hit messages, the collector can preserve the exact received bytes and
still parse the contents for live counters. This design supports later trace
dumping and parser improvement without changing the producer protocol.

Trace dumpers exist in Tcl, Java, C++, and Rust. They can print legacy
stack-depth-style output, parsed hit/log/context lines, raw frame summaries, or
time-filtered subsets. The dumpers support time windows expressed as raw
nanosecond values or RFC3339 UTC timestamps, using the trace header and
timestamp frames to align monotonic file time with wall-clock time.

\section{Testing}

The prototype includes a cross-language regression harness. Each test starts a
collector, runs an instrumented producer or demo application, issues admin
queries, and compares the observed output to a golden transcript. The scenarios
include:

\begin{itemize}
  \item Java server receiving Java probes,
  \item Java server receiving C++ probes,
  \item Java server receiving Rust probes,
  \item Rust server receiving C++ probes,
  \item Java server receiving a JAR-instrumented branch-probe demo,
  \item Java server receiving a fractal/demo workload.
\end{itemize}

These tests matter because the most important property of the system is
interoperability. A producer and collector written in different languages must
agree on byte order, field widths, batching behavior, hit counts, log framing,
stack depth, and admin output. Golden files make that agreement visible.

\section{Evaluation}

\subsection{Latency Benchmarks}

The prototype includes latency benchmark drivers for Java, C++, and Rust. The
benchmark pattern is simple: warm up, execute one million probe hits, measure
elapsed time around the hit calls, and allow the sender thread to flush. An
engineering benchmark note records the following representative measurements:

\begin{center}
\begin{tabular}{llll}
\textbf{Language} & \textbf{Average per hit} & \textbf{Total for 1M hits} & \textbf{Notes} \\
\hline
Java & \(2.6\)--\(2.8\,\mu s\) & about \(2.6\)--\(2.8\,s\) & with stack walking \\
Java & \(0.018\,\mu s\) & \(17.662\,ms\) & constant depth \\
C++ & \(0.071\,\mu s\) & \(71.516\,ms\) & optimized ring path \\
Rust & \(0.115\)--\(0.122\,\mu s\) & about \(115\)--\(122\,ms\) & release build \\
\end{tabular}
\end{center}

These figures should be read as engineering measurements from one prototype
environment, not as a universal benchmark. They are still useful because they show
where the costs are. Java stack walking dominates the Java path. When stack
depth is represented by a cheap counter or constant, Java can be comparable to
native paths for the enqueue portion. C++ and Rust benefit from RAII/drop-based
depth guards and direct control over memory layout.

\subsection{Throughput and Loss Tradeoffs}

The UDP path favors throughput and low perturbation. The runtime can continue
even if a collector is absent. Under overload, drops are possible at the
runtime queue, socket buffer, network, or collector receive buffer. The
implementation mitigates this with batching and large socket buffers, but it
does not pretend UDP is lossless.

This tradeoff is appropriate for many coverage and observability uses. If a
test is checking whether a location was exercised at least once, occasional
duplicate or missing hits may be acceptable when repeated across runs. If a
workflow requires exact event accounting, the same message format can be sent
over TCP or persisted locally before export.

\section{Use Cases}

\subsection{Distributed Coverage}

The most direct use case is coverage across several services. Each service is
assigned an application and instance id. Tests run against the composed system,
and collectors receive hits from every instrumented process. Coverage reports
can then show whether a scenario exercised the expected locations across
language and process boundaries.

\subsection{Branch Coverage for Java Services}

The JAR instrumenter can insert probes at method entries and branch edges. The
probe map records whether a location is \code{METHOD\_ENTRY}, \code{IF\_TRUE},
\code{IF\_FALSE}, \code{CATCH\_ENTRY}, or \code{FINALLY\_ENTRY}. Combining that
map with collector hit counts gives branch coverage without requiring a
language-specific coverage database at runtime.

\subsection{Operational Context Slicing}

Context messages let the observer label a time interval or runtime phase. For
example, a harness can apply \code{scenario=retry-timeout}, run a test, and
withdraw it. Later, the coverage report distinguishes hits observed under that
context from hits observed under \code{default}. This is useful when the same
code is executed by many tests or tenants and the question is which condition
caused a path to run.

\subsection{Trace Replay and Debugging}

Because collectors persist raw frames, a session can be inspected after the
application has exited. Trace dumpers can show hit sequences, log messages, and
context changes. Time filters can restrict inspection to a short interval
around an incident or test phase.

\section{Design Discussion}

\subsection{Why a Small Binary Protocol?}

A small binary protocol is easier to implement consistently across languages
and cheap to produce at high frequency. The 20-byte hit message avoids runtime
allocation for strings, JSON serialization, schema registries, and large trace
objects. The cost is that human-readable source information must be handled by
sidecar metadata such as the branch-probe CSV. That separation is a good trade
for a dense probe system.

\subsection{Why UDP?}

UDP is attractive for high-frequency probes because it decouples application
threads from collector availability. A service can be instrumented even when
the collector is not running. Batching keeps packet count manageable. The
system does not require a connection handshake or backpressure on the common
path. TCP remains useful for modes where loss is unacceptable, and the
collectors already parse both UDP datagrams and TCP streams.

\subsection{Why Context Sets Instead of a Single Context Id?}

A single context label is too weak for real tests and operations. At any
moment, an application may be under several labels: a test scenario, a feature
flag, an experiment, and a tenant class. The collector therefore treats the
active context as a set. Each distinct set receives an id. This keeps hit
records compact while preserving multi-label attribution.

\subsection{Why Multiple Collector Implementations?}

Multiple collector implementations are partly practical and partly validating.
They make it possible to run the system in environments where Java, C++, or
Rust is the natural deployment choice. More importantly, they test whether the
protocol is truly language-neutral. If three collectors and three producer
runtimes can interoperate, accidental assumptions in one implementation are
more likely to be exposed.

\section{Limitations}

The current system has several limitations.

\begin{itemize}
  \item \textbf{UDP loss is possible.} The common path prioritizes low
  overhead. Exact accounting requires TCP or an additional reliability layer.

  \item \textbf{The active context set is global per collector.} This is simple,
  but it may be too coarse when several independent scenarios send to one
  collector at the same time. A future protocol could carry context id or
  request id directly in hit records.

  \item \textbf{Location maps are external to hit records.} This keeps hits
  small but requires careful archival of \code{branch-probes.csv} or equivalent
  metadata for each instrumented build.

  \item \textbf{Java stack depth is expensive when computed by walking the
  stack.} The newer runtime design addresses this with instrumented
  \code{scope\_ENTER}/\code{scope\_EXIT}, but manually instrumented Java code
  must avoid full stack walking in hot paths.

  \item \textbf{The protocol is intentionally minimal.} It does not yet encode
  request ids, parent-child spans, wall-clock timestamps per producer, source
  hashes, or schema versions beyond the trace-file magic.

  \item \textbf{Some tooling is prototype-grade.} The current implementation
  includes demos, experiments, shell scripts, generated build files, and
  several versions of the same idea. The core protocol is clear, but packaging
  and operational polish remain future work.
\end{itemize}

\section{Future Work}

Several extensions follow naturally from the current implementation.

\begin{enumerate}
  \item \textbf{Protocol versioning.} Add an explicit version or capability
  field to support future hit formats while preserving compatibility.

  \item \textbf{Per-hit context ids.} Let producers attach a context id or
  request id to each hit when concurrent context streams share one collector.

  \item \textbf{Reliable mode.} Provide a first-class TCP or local-spool mode
  for tests that require exact hit counts.

  \item \textbf{Trace indexing.} Build indexes over \hittrace\ files so time
  windows, location ids, and context transitions can be queried without a full
  scan.

  \item \textbf{Source-map integration.} Bundle probe maps, build ids, and
  source hashes so reports can be resolved automatically against the correct
  source revision.

  \item \textbf{Unified CLI.} Wrap server startup, instrumentation, trace
  dumping, and coverage report generation in one command-line interface.

  \item \textbf{OpenTelemetry bridge.} Export selected probe events or coverage
  summaries into existing observability systems while preserving the cheap
  internal protocol.
\end{enumerate}

\section{Conclusion}

\tool\ demonstrates that useful distributed coverage can be built from a small
cross-language probe protocol. The design avoids putting analysis work in the
application process. Probes enqueue compact binary hits, background senders
batch them, collectors persist raw history, and admin tools compute
context-aware reports. The prototype validates the idea from several angles:
manual probes, Java bytecode instrumentation, optimized native and JVM
runtimes, Java/C++/Rust collectors, binary trace dumpers, regression tests, and
benchmarks. The result is not a replacement for full tracing infrastructure,
but a focused substrate for asking a valuable question in distributed systems:
which parts of the system actually ran, under which runtime contexts, and with
how little measurement cost?

\begin{thebibliography}{9}

\bibitem{asm}
OW2 Consortium.
\newblock ASM: Java bytecode manipulation and analysis framework.
\newblock \url{https://asm.ow2.io/}.

\bibitem{jacoco}
JaCoCo Project.
\newblock JaCoCo Java Code Coverage Library.
\newblock \url{https://www.jacoco.org/jacoco/}.

\bibitem{opentelemetry}
OpenTelemetry Authors.
\newblock OpenTelemetry specification.
\newblock \url{https://opentelemetry.io/}.

\bibitem{tcl}
Tcl Core Team.
\newblock Tcl scripting language.
\newblock \url{https://www.tcl-lang.org/}.

\bibitem{molt}
Molt Project.
\newblock Molt Tcl-like interpreter for Rust.
\newblock \url{https://github.com/wduquette/molt}.

\end{thebibliography}

\end{document}
